\section{El camino del código abierto: Tianshou y tuixue.online}

\subsection{Tianshou (天授): un marco de trabajo de RL hecho a mano en dos semanas}

A principios de 2020, Weng Jiayi (翁家翌) tomó una decisión que cambiaría el rumbo de su carrera: escribir desde cero un marco de trabajo de aprendizaje por refuerzo.

El motivo era sencillo: había escrito mucho código de experimentos de RL y quería integrarlo para que corriera mejor. Primero revisó RLlib, dentro de Ray, y tras un mes descubrió que era demasiado complejo: cientos de miles de líneas de código, demasiadas abstracciones, sin ninguna idea de por dónde modificarlo. Así que decidió \textbf{tirarlo todo y empezar de cero, construyendo a mano un marco de trabajo completamente nuevo}.

\begin{importantbox}{La filosofía de diseño de Tianshou: la consistencia (Consistency)}
Weng Jiayi considera que la característica más importante de un buen proyecto es la \textbf{consistencia}. Los proyectos en los que colaboran muchas personas tienden a degradarse, porque nadie sabe bien qué escribió el otro; los supuestos no se transmiten a tiempo, lo que provoca que el código se copie y pegue y se infle. Tianshou tuvo éxito porque una sola persona lo diseñó e implementó de principio a fin, lo que garantizó la consistencia del conjunto.
\end{importantbox}

La primera versión de Tianshou se completó en solo \textbf{dos semanas}. Weng Jiayi explica que, si las abstracciones están bien planteadas, implementar un algoritmo puede tomar menos de veinte líneas de código. Frente a las cientos de miles de líneas de RLlib, la clave de Tianshou está en un diseño de alto nivel reducido al mínimo: si el usuario quiere modificar alguna función, el diseño ya especifica de antemano «solo se puede modificar aquí».

Lo que Tianshou hizo bien en el fondo fue \textbf{captar la necesidad de los usuarios}: en ese momento, los investigadores de RL necesitaban un marco de trabajo fácil de usar, fácil de modificar y con poco código, que pudieran tomar directamente y, con solo estudiarlo un poco, supieran dónde modificar. Tianshou satisfizo exactamente esa necesidad.

\subsubsection{El desarrollo posterior de Tianshou y su degradación}

Tianshou pasó después a ser un proyecto mantenido por la comunidad de código abierto. Cuando Weng Jiayi entró a trabajar en OpenAI, ya no tuvo tiempo de seguir manteniéndolo, así que transfirió el mantenimiento a la comunidad. La regla que estableció fue: mientras haya una sola persona que tome las decisiones finales, se puede mantener la consistency.

Cinco años después, al mirar atrás, Weng Jiayi reconoce que Tianshou se «degradó un poco», porque su context y el context de quien lo sucedió no coincidían del todo, y el sucesor reescribió parte del código, lo que hizo que el conjunto fuera menos consistent. Pero considera que, a largo plazo, esto es aceptable.

Esta experiencia también lo llevó a tomar, en agosto de 2022, una decisión importante: dejar de desarrollar Tianshou de manera gradual. La razón fue que, tras entrar a OpenAI, se dio cuenta de que existía una brecha enorme entre los toy benchmarks a los que apuntaba Tianshou —como Atari y MuJoCo— y el RL que la industria realmente necesitaba (por ejemplo, la alineación de LLM). «Debería invertir más tiempo en cosas más significativas».

\begin{warningbox}{Sobre la decisión de «no publicar papers»}
Weng Jiayi afirma con claridad que no hizo Tianshou para publicar artículos: «Creo que publicar papers no tiene ningún sentido». En ese momento ya tenía artículos publicados y resultados en competencias, suficientes para sus solicitudes. Lo que quería era un proyecto de código abierto real, con más de tres dígitos de stars, un proyecto de GitHub «como es debido», según el sistema de evaluación de su asesor.
\end{warningbox}

\subsection{tuixue.online: un sistema de consulta de visas durante la pandemia}

Durante la pandemia de 2020, los consulados de Estados Unidos cerraron y los tiempos de espera de visa eran sumamente inciertos. El propio Weng Jiayi también estaba esperando su visa para ir a CMU, así que escribió un rastreador de tiempos de visa y lo publicó como código abierto en forma de un sistema de consulta gratuito: tuixue.online.

La primera versión de este proyecto ni siquiera necesitaba mucha tecnología: se actualizaba a mano una vez de día y una vez de noche. Pero incluso así, tan rudimentario, la demanda ya era enorme. Después de automatizarlo, el número acumulado de clics superó el millón (y más adelante quizá llegó a los diez millones). Cuando terminó la pandemia, los consulados actualizaron sus sitios web, el rastreador original dejó de funcionar y Weng Jiayi tampoco tuvo tiempo de reescribirlo, así que el proyecto cumplió su cometido histórico.

La conclusión de Weng Jiayi sobre este proyecto es: «La tecnología no importa; lo que importa es captar la necesidad». Ni siquiera se necesita tecnología avanzada: incluso una actualización manual, si satisface el punto de dolor real de los usuarios, tiene un valor enorme.

\subsubsection{El código abierto como una forma de caridad}

Weng Jiayi define reiteradamente sus proyectos de código abierto como «caridad». El presentador señala que, en realidad, ni Tianshou ni tuixue.online son proyectos utilitarios: cuando hizo Tianshou, sus solicitudes ya habían terminado, y tuixue es completamente gratuito. Weng Jiayi reconoce que tiene un «impulso interno muy fuerte»: el de crear algo que él considera útil y compartirlo con todos.

Compartió el origen de este impulso: en el último año de preparatoria le surgió de pronto una idea: «Si la vida es un juego, entonces la puntuación final del juego es el número de personas que recuerdan tu nombre en el instante en que mueres». Esta «teoría de la vida como juego» lo ha impulsado desde entonces a buscar impact, y sigue haciéndolo hasta hoy.

El presentador insistió: «¿No sientes que tú también terminas empujado por este criterio?». La respuesta de Weng Jiayi resulta interesante: «Por ahora no. Si me doy cuenta de que pasa eso, puedo cambiarlo; puedo cambiar mi propio criterio de evaluación». No es «esclavo» de ese criterio: aunque durante mucho tiempo en OpenAI no haya tenido ningún proyecto nuevo de código abierto, tampoco le preocupa, porque «el OpenAI model ya es el mejor (proyecto de código abierto)».

\begin{importantbox}{El código como caridad}
Weng Jiayi define tanto Tianshou como tuixue.online como \textbf{proyectos de caridad} (non-profit). Dice: «Hacer proyectos de caridad me da una satisfacción enorme. Comparado con el dinero, el impact me satisface más». Esta idea de «el código como caridad» viene de una «regla de puntuación del juego de la vida» de la que se dio cuenta de pronto en la preparatoria: \textbf{cuantas más personas recuerden tu nombre en el instante en que mueres, más alta es la puntuación}.
\end{importantbox}

\subsection{Resumen de la sección}

Tianshou y tuixue.online comparten el mismo patrón de creación: \textbf{tener una necesidad propia → no existir en el mercado una herramienta buena → construir una a mano → publicarla como código abierto y gratuita para todos}. Ninguno de los dos proyectos es utilitario: Tianshou no se hizo para publicar artículos, tuixue no se hizo para ganar dinero. Pero fueron precisamente estos proyectos «no utilitarios» los que le trajeron a Weng Jiayi el mayor impact y las mayores oportunidades profesionales. La tecnología no importa; lo que importa es \textbf{captar la necesidad}.

